% ============================================================
\chapter{Bases de Donn\'ees NoSQL}
\label{ch:nosql}
% ============================================================

\begin{intuition}[title=\textbf{Id\'ee centrale}]
Les bases NoSQL (\emph{Not Only SQL}) abandonnent le mod\`ele relationnel
strict pour gagner en scalabilit\'e horizontale, en flexibilit\'e de
sch\'ema et en performances sur des cas d'usage sp\'ecifiques. Quatre
grandes familles existent~: cl\'e-valeur, document, colonne et graphe.
\end{intuition}

% ------------------------------------------------------------
\section{Motivation et contexte}
% ------------------------------------------------------------

Les bases relationnelles atteignent leurs limites face \`a~:
\begin{itemize}
  \item des volumes massifs (p\'etaoctets),
  \item des donn\'ees semi-structur\'ees ou non structur\'ees,
  \item des besoins de haute disponibilit\'e et de distribution g\'eographique,
  \item des sch\'emas \'evolutifs (d\'eveloppement agile).
\end{itemize}

\begin{remarque}
NoSQL ne signifie pas \og anti-SQL \fg mais \og Not Only SQL \fg.
Beaucoup de bases NoSQL proposent d\'esormais un langage de requ\^ete
d\'eclaratif (CQL pour Cassandra, MQL pour MongoDB, Cypher pour Neo4j).
\end{remarque}

% ------------------------------------------------------------
\section{Th\'eor\`eme CAP}
% ------------------------------------------------------------

\begin{theoreme}[Th\'eor\`eme CAP (Brewer, 2000)]
\index{th\'eor\`eme CAP}
Un syst\`eme distribu\'e ne peut garantir simultan\'ement que deux des
trois propri\'et\'es suivantes~:
\begin{enumerate}
  \item \textbf{Consistency} (coh\'erence)~: toute lecture renvoie la
        derni\`ere \'ecriture.
  \item \textbf{Availability} (disponibilit\'e)~: toute requ\^ete re\c{c}oit
        une r\'eponse (non n\'ecessairement la plus r\'ecente).
  \item \textbf{Partition tolerance} (tol\'erance aux partitions)~: le
        syst\`eme fonctionne malgr\'e des coupures r\'eseau.
\end{enumerate}
\end{theoreme}

\begin{attention}
En pratique, les partitions r\'eseau sont in\'evitables dans un syst\`eme
distribu\'e. Le vrai choix est donc entre \textbf{CP} (coh\'erence forte,
ex.~: HBase, MongoDB avec \texttt{readConcern: majority}) et \textbf{AP}
(disponibilit\'e, ex.~: Cassandra, DynamoDB en mode \'eventual).
\end{attention}

\begin{definition}[Coh\'erence \'eventuelle]
\index{coh\'erence \'eventuelle}
Un syst\`eme \`a \textbf{coh\'erence \'eventuelle} (\emph{eventual consistency})
garantit que, en l'absence de nouvelles \'ecritures, toutes les r\'epliques
convergent vers la m\^eme valeur apr\`es un d\'elai fini.
\end{definition}

% ------------------------------------------------------------
\section{Bases cl\'e-valeur}
% ------------------------------------------------------------

\begin{definition}[Base cl\'e-valeur]
\index{base cl\'e-valeur}
Une \textbf{base cl\'e-valeur} stocke des paires $(k, v)$ o\`u $k$ est
une cl\'e unique et $v$ une valeur opaque (cha\^ine, blob, objet s\'erialis\'e).
L'acc\`es se fait uniquement par cl\'e~: \texttt{GET}, \texttt{SET},
\texttt{DELETE}.
\end{definition}

\begin{exemple}
Redis~: cache en m\'emoire et structure de donn\'ees~:
\begin{minted}{python}
import redis

r = redis.Redis(host='localhost', port=6379, db=0)
r.set('user:1001:name', 'Alice')
r.set('user:1001:email', 'alice@example.com')
r.expire('user:1001:name', 3600)  # TTL de 1 heure

name = r.get('user:1001:name')  # b'Alice'

# Structures avancees
r.hset('user:1002', mapping={'name': 'Bob', 'age': '25'})
r.lpush('queue:jobs', 'task_42')
r.zadd('leaderboard', {'Alice': 1500, 'Bob': 1200})
\end{minted}
\end{exemple}

\begin{remarque}
Cas d'usage typiques~: cache applicatif, sessions utilisateur,
compteurs temps r\'eel, files d'attente de messages.
\end{remarque}

% ------------------------------------------------------------
\section{Bases document}
% ------------------------------------------------------------

\begin{definition}[Base document]
\index{base document}
Une \textbf{base document} stocke des documents semi-structur\'es
(JSON, BSON) regroup\'es en \emph{collections}. Chaque document a un
sch\'ema flexible~: deux documents de la m\^eme collection peuvent avoir
des champs diff\'erents.
\end{definition}

\begin{exemple}
MongoDB~:
\begin{minted}{python}
from pymongo import MongoClient

client = MongoClient('mongodb://localhost:27017/')
db = client['universite']
col = db['etudiants']

# Insertion
col.insert_one({
    'nom': 'Dupont',
    'prenom': 'Alice',
    'notes': [
        {'matiere': 'BDD', 'note': 16},
        {'matiere': 'Algo', 'note': 14}
    ],
    'adresse': {'ville': 'Paris', 'cp': '75005'}
})

# Requete avec filtre et projection
resultats = col.find(
    {'notes.note': {'$gte': 15}},
    {'nom': 1, 'prenom': 1, '_id': 0}
)

# Agregation (equivalent GROUP BY)
pipeline = [
    {'$unwind': '$notes'},
    {'$group': {'_id': '$nom', 'moyenne': {'$avg': '$notes.note'}}}
]
col.aggregate(pipeline)
\end{minted}
\end{exemple}

\begin{remarque}
MongoDB offre des index secondaires, des transactions multi-documents
(depuis la version 4.0) et un langage d'agr\'egation riche. Il convient
aux donn\'ees \`a sch\'ema variable (catalogues produits, CMS, IoT).
\end{remarque}

% ------------------------------------------------------------
\section{Bases orient\'ees colonnes}
% ------------------------------------------------------------

\begin{definition}[Base orient\'ee colonnes]
\index{base orient\'ee colonnes}
Une \textbf{base orient\'ee colonnes} (ou \emph{column-family store})
organise les donn\'ees par familles de colonnes plut\^ot que par lignes.
Chaque ligne est identifi\'ee par une cl\'e de partition et peut avoir un
nombre variable de colonnes.
\end{definition}

\begin{exemple}
Apache Cassandra~:
\begin{minted}{python}
-- CQL (Cassandra Query Language)
CREATE KEYSPACE universite WITH replication = {
    'class': 'SimpleStrategy', 'replication_factor': 3
};

CREATE TABLE universite.notes (
    etudiant_id UUID,
    matiere TEXT,
    semestre INT,
    note FLOAT,
    PRIMARY KEY ((etudiant_id), matiere, semestre)
) WITH CLUSTERING ORDER BY (matiere ASC, semestre DESC);

-- La cle de partition (etudiant_id) determine le noeud de stockage
-- Les colonnes de clustering ordonnent les donnees au sein de la partition
\end{minted}
\end{exemple}

\begin{remarque}
Cassandra excelle pour les \'ecritures massives distribu\'ees (logs,
s\'eries temporelles, IoT). La mod\'elisation est pilot\'ee par les
requ\^etes~: on con\c{c}oit les tables en fonction des \texttt{SELECT}
que l'on souhaite ex\'ecuter.
\end{remarque}

% ------------------------------------------------------------
\section{Bases de donn\'ees orient\'ees graphe}
% ------------------------------------------------------------

\begin{definition}[Base de donn\'ees graphe]
\index{base de donn\'ees graphe}
Une \textbf{base de donn\'ees graphe} stocke des \textbf{n\oe uds}
(entit\'es) et des \textbf{ar\^etes} (relations) comme citoyens de
premi\`ere classe. Les travers\'ees de graphe sont des op\'erations natives
en $O(1)$ par relation (ind\'ependamment de la taille du graphe).
\end{definition}

\begin{exemple}
Neo4j et le langage Cypher~:
\begin{minted}{python}
// Creer des noeuds et des relations
CREATE (a:Personne {nom: 'Alice', age: 30})
CREATE (b:Personne {nom: 'Bob', age: 25})
CREATE (a)-[:AMIS_AVEC {depuis: 2020}]->(b)
CREATE (a)-[:TRAVAILLE_A]->(:Entreprise {nom: 'TechCorp'})

// Trouver les amis d'amis (distance 2)
MATCH (p:Personne {nom: 'Alice'})-[:AMIS_AVEC*2]-(fof)
WHERE fof <> p
RETURN DISTINCT fof.nom

// Plus court chemin
MATCH path = shortestPath(
    (a:Personne {nom: 'Alice'})-[*..6]-(b:Personne {nom: 'Charlie'})
)
RETURN path
\end{minted}
\end{exemple}

\begin{remarque}
Les bases graphe sont id\'eales pour les r\'eseaux sociaux, la d\'etection
de fraude, les syst\`emes de recommandation et la gestion des connaissances
(knowledge graphs). Elles surpassent les bases relationnelles d\`es que les
requ\^etes impliquent des jointures r\'ecursives profondes.
\end{remarque}

% ------------------------------------------------------------
\section{Comparaison et choix}
% ------------------------------------------------------------

\begin{center}
\begin{tabular}{lcccc}
\hline
& \textbf{Cl\'e-valeur} & \textbf{Document} & \textbf{Colonne} & \textbf{Graphe} \\
\hline
Exemple    & Redis    & MongoDB    & Cassandra & Neo4j \\
Sch\'ema   & Aucun    & Flexible   & Semi-fixe & Flexible \\
Requ\^etes & Par cl\'e & Riches     & Par partition & Travers\'ees \\
Scalabilit\'e & Tr\`es haute & Haute & Tr\`es haute & Mod\'er\'ee \\
Cas typique & Cache & CMS, catalogue & S\'eries temp. & R\'eseau social \\
\hline
\end{tabular}
\end{center}

\begin{proposition}[Crit\`eres de choix]
Le choix d'un type de base NoSQL d\'epend de~:
\begin{enumerate}
  \item la \textbf{structure} des donn\'ees (sch\'ema fixe vs.\ flexible),
  \item les \textbf{patterns d'acc\`es} (cl\'e, requ\^etes complexes, travers\'ees),
  \item les exigences de \textbf{coh\'erence} vs.\ \textbf{disponibilit\'e},
  \item le \textbf{volume} et la \textbf{v\'elocit\'e} des donn\'ees.
\end{enumerate}
\end{proposition}

% ------------------------------------------------------------
\section{Formules cl\'es}
% ------------------------------------------------------------

\begin{formules}
\begin{itemize}
  \item \textbf{CAP}~: Consistency + Availability + Partition tolerance --- au plus 2 sur 3
  \item \textbf{Cl\'e-valeur}~: acc\`es $O(1)$ par cl\'e, aucune requ\^ete sur la valeur
  \item \textbf{Document}~: sch\'ema flexible, index secondaires, agr\'egation
  \item \textbf{Colonne}~: mod\'elisation par requ\^ete, \'ecritures distribu\'ees massives
  \item \textbf{Graphe}~: travers\'ees en $O(k)$ o\`u $k$ = nombre d'ar\^etes travers\'ees
\end{itemize}
\end{formules}

% ------------------------------------------------------------
\section{Exercices}
% ------------------------------------------------------------

\begin{exercice}
Pour chacun des cas d'usage suivants, recommander un type de base NoSQL
et justifier~: (a) cache de sessions web, (b) catalogue de produits avec
attributs variables, (c) suivi de capteurs IoT \`a 100\,000 \'ev\'enements/s,
(d) d\'etection de communaut\'es dans un r\'eseau social.
\end{exercice}

\begin{exercice}
\'Ecrire un script Python qui ins\`ere 10\,000 documents dans MongoDB,
cr\'ee un index sur un champ, puis compare le temps d'une requ\^ete avec
et sans index.
\end{exercice}

\begin{exercice}
Mod\'eliser un syst\`eme de recommandation de films dans Neo4j. Cr\'eer
les n\oe uds (utilisateurs, films, genres) et les relations (A\_VU,
A\_NOTE, EST\_DU\_GENRE). \'Ecrire une requ\^ete Cypher pour recommander
des films \`a un utilisateur bas\'e sur les pr\'ef\'erences de ses amis.
\end{exercice}

\begin{exercice}
Concevoir le sch\'ema Cassandra pour une application de messagerie
(messages entre utilisateurs). La requ\^ete principale est \og afficher
les $N$ derniers messages d'une conversation \fg. Justifier le choix
de la cl\'e de partition et des colonnes de clustering.
\end{exercice}

\begin{exercice}
Discuter les implications du th\'eor\`eme CAP pour une application
bancaire distribu\'ee sur plusieurs continents. Quel compromis
recommanderiez-vous~? Pourquoi~?
\end{exercice}
